% !TEX root = ../main.tex
\chapter{Bridge Boundaries to the External World}
\label{ch:bridge_boundaries}

\Lucid's engine is a closed system. Inside \code{lucid/} there are
no imports of \code{numpy}, no calls into \code{scipy}, no
dependence on any third-party numerical library other than \MLX\
and \Accelerate, both of which are themselves wrapped by \Lucid's
backends. The closure is enforced by Hard Rule H4
(\appref{app:hard_rules}), which is in turn enforced by a static
check that scans every Python file in \code{lucid/} for forbidden
imports.

But the engine is not a sealed system. It must accept data from the
outside (an image dataset arriving as numpy arrays, a checkpoint
loaded from disk), and it must hand data to the outside (a tensor
converted to numpy for plotting, a model state serialised for
distribution). The reconciliation is a set of six explicitly named
\emph{bridge boundaries}: six places in \code{lucid/} where
external libraries are permitted, each justified by a specific
interoperability need and each documented in this chapter.

Any place outside these six where a forbidden import appears is a
hard-rule violation that fails CI. This chapter exists to enumerate
the boundaries, justify each one, describe what crosses each
boundary, and document the discipline that keeps them as the only
contact points with the outside world.

The boundaries, in order:

\begin{enumerate}
  \item \code{lucid/\_factories/converters.py}: external array
    ingest --- \code{tensor}, \code{from\_numpy},
    \code{from\_dlpack}, \code{to\_dlpack}.
  \item \code{lucid/\_tensor/tensor.py}: tensor-side egress ---
    \code{Tensor.numpy()}, \code{Tensor.\_\_dlpack\_\_},
    \code{Tensor.\_to\_impl}.
  \item \code{lucid/\_tensor/\_repr.py}: display only ---
    \code{\_\_repr\_\_} and \code{\_\_str\_\_}.
  \item \code{lucid/\_types.py}: typing protocol ---
    \code{TYPE\_CHECKING}-only imports.
  \item \code{lucid/serialization/} and
    \code{lucid/optim/\{optimizer,lbfgs\}.py}: checkpoint
    serialisation --- \code{state\_dict}, \code{load\_state\_dict}.
  \item \code{lucid/utils/data/dataloader.py}: external data
    ingest --- \code{DataLoader}'s collate path.
\end{enumerate}

\tabref{tab:bridge_boundaries:summary} gives the same information
in a quick-reference form.

\begin{table}[t]
\centering
\footnotesize
\setlength{\tabcolsep}{3pt}
\begin{adjustbox}{max width=\textwidth, center}
\begin{tabular}{rllll}
\toprule
\# & File & Direction & External lib & Use case \\
\midrule
1 & \code{\_factories/converters.py} & ingest  & numpy, dlpack & \code{tensor()}, \code{from\_numpy} \\
2 & \code{\_tensor/tensor.py}        & egress  & numpy, dlpack & \code{.numpy()}, \code{\_\_dlpack\_\_} \\
3 & \code{\_tensor/\_repr.py}        & display & numpy         & \code{\_\_repr\_\_} formatting \\
4 & \code{\_types.py}                & typing  & numpy (TYPE\_CHECKING) & \code{NDArrayLike} alias \\
5 & \code{serialization/}, \code{optim/optimizer.py} & checkpoint & numpy, safetensors & \code{state\_dict} I/O \\
  & \code{optim/lbfgs.py} & checkpoint & numpy & LBFGS history \\
6 & \code{utils/data/dataloader.py}  & ingest  & numpy         & default \code{collate\_fn} \\
\bottomrule
\end{tabular}
\end{adjustbox}
\caption[The six bridge boundaries.]{The exhaustive list of places
in \code{lucid/} where an external numerical library may be
imported. Any \texttt{import numpy} (or scipy, torch, jax, etc.)
outside these six files is a Hard Rule H4 violation rejected by
\texttt{tools/check\_bridges.py} in CI.}
\label{tab:bridge_boundaries:summary}
\end{table}

\secref{sec:bridge_boundaries:rule} states the rule.
Sections~\ref{sec:bridge_boundaries:b1}--\ref{sec:bridge_boundaries:b6}
treat each boundary in turn.
\secref{sec:bridge_boundaries:enforcement} describes how the rule
is enforced. \secref{sec:bridge_boundaries:adding} describes the
policy for adding (or rejecting) a seventh boundary.

\section{Hard Rule H4 Restated}
\label{sec:bridge_boundaries:rule}

The verbatim text of H4:

\begin{lucidrule}
Lucid internal behaviour (\code{lucid/}'s composite, tensor, nn,
optim, autograd, ops, linalg, fft, signal, special, distributions,
einops, amp, profiler) must not depend on any external library.
NumPy, SciPy, or any other external package may not be imported.
Only the Lucid C++ engine and the Python standard library are
permitted. The sole exception is the six bridge boundaries
enumerated below.
\end{lucidrule}

The motivation is twofold:

\begin{enumerate}
  \item \textbf{Numerical hermeticism.} An external numerical
    library injected into a computational path can subtly change
    the dtype, the layout, or the numerical result of an operation.
    Such a change is invisible at the import line and emerges only
    when a parity test against the reference framework fails.
    Keeping the engine numpy-free means every numerical decision is
    made by \Lucid\ code that can be audited and changed under our
    own discipline.
  \item \textbf{Dependency surface.} Every external import is a
    place where a future \code{pip install} can fail, a version
    constraint can drift, or a security patch must be tracked.
    The six bridge boundaries are explicitly version-pinned and
    quarantined; the rest of the engine is unaffected by changes
    in the external ecosystem.
\end{enumerate}

The rule was established in early \Lucid\ 3.x and tightened in 3.0.2
when the framework completed its numpy-independence pass (vault:
\code{retro-3-0-2-numpy-free-standalone}). The six boundaries
listed here are the residue of that pass.

\section{Boundary 1: External Array Ingest}
\label{sec:bridge_boundaries:b1}

\code{lucid/\_factories/converters.py} is the file that constructs
\Lucid\ tensors from external array-like objects. It imports
\code{numpy} unconditionally (numpy is a required dependency) and
\code{torch}-style DLPack producers conditionally (DLPack is the
standard cross-framework tensor protocol).

The functions exposed:

\begin{itemize}
  \item \code{lucid.tensor(data, dtype=None, device=None)}:
    construct a tensor from a Python scalar, list, numpy array, or
    array-like. The function inspects the input type and routes
    accordingly; numpy arrays go through \code{from\_numpy}, lists
    are converted via numpy first, scalars become 0-d tensors.
  \item \code{lucid.from\_numpy(arr)}: zero-copy construction from
    a contiguous numpy array. The numpy buffer is wrapped in a
    \code{Storage} with \code{provider = External} and the numpy
    array is held alive for the storage's lifetime through a
    closure deleter.
  \item \code{lucid.from\_dlpack(capsule)}: construction from a
    DLPack capsule \cite{dlpack}. Same zero-copy mechanic; the
    capsule's deleter is invoked when the storage is freed.
  \item \code{lucid.to\_dlpack(tensor)}: production of a DLPack
    capsule from a \Lucid\ tensor, for interop with other DLPack
    consumers.
\end{itemize}

\subsection{Why DLPack is routed through numpy}
\label{ssec:bridge_boundaries:dlpack_numpy}

The DLPack path in \Lucid\ goes through numpy as an intermediate.
Specifically, \code{from\_dlpack} first calls
\code{numpy.from\_dlpack(capsule)} and then \code{from\_numpy} on
the resulting numpy array. The alternative would have been to
construct a \Lucid\ tensor directly from the DLPack capsule, which
would avoid the numpy round-trip.

The alternative was implemented and abandoned. The vault note
\code{arch-dlpack-via-numpy} records the reasoning: the direct
DLPack path was 350 lines of C code that handled dtype mapping,
stride translation, and capsule lifetime; the indirect path through
numpy is 5 lines of Python; the performance difference is
unmeasurable (DLPack is a metadata transfer, no data motion) and
numpy's DLPack implementation is well-tested. The maintenance cost
of the direct path was deemed not worth the marginal architectural
purity.

\section{Boundary 2: Tensor-Side Egress}
\label{sec:bridge_boundaries:b2}

\code{lucid/\_tensor/tensor.py} contains the \code{Tensor} class
definition. The class has three methods that touch numpy or
DLPack:

\begin{itemize}
  \item \code{Tensor.numpy()}: return a numpy view of the tensor's
    data. The tensor must be on \code{device=cpu}; for a metal
    tensor, the user must first call \code{.to('cpu')}.
  \item \code{Tensor.\_\_dlpack\_\_()}: the dunder that
    third-party DLPack consumers call to obtain a capsule.
  \item \code{Tensor.\_to\_impl(device)}: internal device-move
    implementation that, in one specific case, uses numpy as a
    fallback for an unusual dtype combination.
\end{itemize}

Of these, \code{.numpy()} is by far the most-used. The
implementation is:

\begin{lstlisting}[style=lucid,language=LucidPython]
def numpy(self):
    if self.device.type != 'cpu':
        raise LucidError("call .to('cpu') before .numpy()")
    if self.requires_grad:
        raise LucidError("call .detach() before .numpy()")
    return _C_engine.tensor_to_numpy(self._impl)
\end{lstlisting}

The two preconditions (\code{.to('cpu')}, \code{.detach()}) are
explicit: \Lucid\ does not silently move to CPU and does not
silently detach. The reasoning is the same as for cross-device
binary ops (\chref{ch:backend_dispatch}): the user should be
intentional about device moves and about breaking the autograd
graph.

The \code{tensor\_to\_numpy} bound C++ function constructs a numpy
array whose buffer points at the \code{TensorImpl}'s storage and
holds a reference to the \code{TensorImpl} through a numpy capsule
deleter. The conversion is zero-copy: the numpy array reads the
same physical bytes as the tensor.

\subsection{Why repr is a separate boundary}
\label{ssec:bridge_boundaries:repr_separate}

A natural question: why is the tensor's \code{\_\_repr\_\_} a
separate boundary (boundary 3) rather than part of
\code{tensor.py}? The answer is organisational: \code{tensor.py}
must import numpy for the egress functions and we want to keep its
internal logic away from display code. \code{\_repr.py} is a
single-purpose file that does only formatting, which keeps the
diff for any repr-related change small and reviewable.

\section{Boundary 3: Display Only}
\label{sec:bridge_boundaries:b3}

\code{lucid/\_tensor/\_repr.py} is the file that implements
\code{Tensor.\_\_repr\_\_} and \code{Tensor.\_\_str\_\_}. It
imports numpy because numpy's pretty-printing
(\code{numpy.array2string}) is the well-tested standard for
displaying multi-dimensional arrays in Python. \Lucid's repr
delegates to it:

\begin{lstlisting}[style=lucid,language=LucidPython]
def _format_tensor(t):
    """Construct a repr string for a Tensor."""
    arr = numpy.asarray(_C_engine.tensor_to_numpy(t._impl))
    body = numpy.array2string(arr, precision=4, suppress_small=True)
    return f"Tensor({body}, dtype={t.dtype}, device={t.device})"
\end{lstlisting}

The function is invoked only by the Python interpreter when the
user asks for a repr (typing the variable name at a REPL prompt,
calling \code{print}, embedding in an f-string). It does not run
during training and contributes no runtime overhead.

\subsection{The display carve-out is read-only}
\label{ssec:bridge_boundaries:repr_readonly}

The repr code reads tensor data and returns a string; it does not
modify the tensor, does not produce a new tensor, and does not
affect any state outside the function. This makes it the least
risky of the six boundaries: even a bug in the repr code (typing
wrong dtype, getting precision wrong) does not corrupt
computation. The boundary is enumerated for completeness rather
than for any correctness concern.

\section{Boundary 4: Typing Protocol}
\label{sec:bridge_boundaries:b4}

\code{lucid/\_types.py} contains the type aliases that the rest of
\Lucid\ uses for static type checking. It imports numpy under a
\code{TYPE\_CHECKING} guard so that the import has no runtime cost:

\begin{lstlisting}[style=lucid,language=LucidPython]
from typing import TYPE_CHECKING

if TYPE_CHECKING:
    import numpy as _np
    NDArrayLike = _np.ndarray
else:
    NDArrayLike = object
\end{lstlisting}

The pattern: at runtime, \code{NDArrayLike} is just \code{object}
(no numpy import); at type-check time, \code{mypy} sees the real
\code{numpy.ndarray} alias. This satisfies the type system without
adding a runtime dependency.

The boundary is named because, technically, the file does import
numpy. The import is dead at runtime but live at lint time. We list
it explicitly to forestall the question of why \code{\_types.py}
contains the word \code{numpy}.

\subsection{Why TYPE\_CHECKING and not a string annotation}
\label{ssec:bridge_boundaries:typing_string}

The alternative would have been to use string annotations:
\code{NDArrayLike = "numpy.ndarray"}. \Lucid\ rejects this pattern
per Hard Rule H7 (no string type hints), which is enforced by mypy
\code{strict} configuration. The \code{TYPE\_CHECKING} block is
the H7-compliant equivalent.

\section{Boundary 5: Checkpoint Serialisation}
\label{sec:bridge_boundaries:b5}

\code{lucid/serialization/} and two files in \code{lucid/optim/}
(\code{optimizer.py} and \code{lbfgs.py}) form the fifth boundary.
The serialisation subsystem reads and writes \code{state\_dict}
files, which by convention contain numpy arrays.

\subsection{The state\_dict format}
\label{ssec:bridge_boundaries:state_dict}

A \code{state\_dict} is a Python dict mapping string keys
(parameter names) to tensor values. \Lucid's \code{save} function
serialises the dict using Python's \code{pickle} format, with each
\Lucid\ tensor converted to a numpy array on the way out. The
inverse, \code{load}, reads the pickle and converts numpy arrays
back to \Lucid\ tensors.

The numpy intermediate is necessary because the pickle format does
not know how to represent a \Lucid\ tensor by itself. We could have
added a custom pickle reducer, but the user-facing benefit is
zero: external tools that read checkpoints (e.g.\ for inference
deployment in non-\Lucid\ environments) expect numpy arrays, and
exposing \Lucid\ tensors would have required the reader to depend
on \Lucid.

\subsection{The safetensors path}
\label{ssec:bridge_boundaries:safetensors}

\code{lucid.serialization} also supports the safetensors format
\cite{safetensors}, which is a tensor-format-only alternative to
pickle that is safer (no arbitrary code execution during load) and
faster (memory-mapped read). The safetensors path uses the
\code{safetensors} Python library, which is itself numpy-aware. The boundary is therefore
double: numpy for the tensor representation, safetensors for the
format. Both are quarantined to \code{lucid/serialization/}.

\subsection{Why lbfgs is part of this boundary}
\label{ssec:bridge_boundaries:lbfgs}

\code{lucid/optim/lbfgs.py} is the only optimiser that maintains
state that does not fit the standard \code{state\_dict} schema
(LBFGS maintains a history of past gradients and parameter steps).
Its \code{state\_dict} / \code{load\_state\_dict} are non-trivial
and touch the numpy boundary directly rather than going through
the generic serialisation path. The optimiser base class
(\code{lucid/optim/optimizer.py}) likewise has a small numpy touch
for parameter group serialisation.

This was a deliberate trade-off: factoring the LBFGS state out into
the generic path would have required generalising the schema in a
way that benefits only LBFGS, and the cost of two extra files in
the bridge enumeration was lower than the cost of the
generalisation.

\section{Boundary 6: External Data Ingest}
\label{sec:bridge_boundaries:b6}

\code{lucid/utils/data/dataloader.py} is the user-facing data
loading utility. Its \code{DataLoader} class accepts a
\code{Dataset} that produces samples, batches them, optionally
collates them, and yields a stream of batched tensors.

The boundary appears in the collate path, which is where samples
become tensors. The default \code{collate\_fn} accepts numpy arrays
and lists of numpy arrays and turns them into \Lucid\ tensors via
\code{lucid.from\_numpy}. It also accepts Python scalars, lists,
and dicts of any of the above.

\subsection{Why the data path is its own boundary}
\label{ssec:bridge_boundaries:why_data}

The data ingest path is the single largest source of external data
volume in a typical training workload --- every image, every text
token, every audio sample passes through it. Routing this through
the existing \code{lucid.tensor} or \code{lucid.from\_numpy}
boundary would have been correct but inconvenient: the
\code{DataLoader} class needs to know about collation, batching,
multi-process worker coordination, and other concerns that
\code{from\_numpy} does not.

The boundary therefore exists for organisational reasons: the data
loading code is concentrated in one place, and that place is
allowed to import numpy. The user can write a custom
\code{collate\_fn} that imports anything; \Lucid's contract is that
the default collate works for numpy-formatted samples.

\section{Enforcement}
\label{sec:bridge_boundaries:enforcement}

The H4 rule is enforced by a static check in
\code{tools/check\_bridges.py}. The check walks every \code{.py}
file under \code{lucid/}, parses imports, and verifies that no
file outside the six named files imports any of the forbidden
modules (numpy, scipy, torch, jax, tensorflow, and a handful of
others).

The check runs as part of the pre-commit hook and as part of CI. A
violation fails the build and the contributor is asked to either
(a) move the new dependency into a bridge file or (b) eliminate
the dependency.

\subsection{The bridge whitelist}
\label{ssec:bridge_boundaries:whitelist}

The whitelist is hard-coded in \code{tools/check\_bridges.py}:

\begin{lstlisting}[style=lucid,language=LucidPython]
BRIDGE_FILES = {
    "lucid/_factories/converters.py",
    "lucid/_tensor/tensor.py",
    "lucid/_tensor/_repr.py",
    "lucid/_types.py",
    "lucid/serialization/__init__.py",
    "lucid/optim/optimizer.py",
    "lucid/optim/lbfgs.py",
    "lucid/utils/data/dataloader.py",
}

FORBIDDEN_OUTSIDE_BRIDGE = {
    "numpy", "scipy", "torch", "jax", "tensorflow", "pandas",
    "sklearn", "xarray", "cupy",
}
\end{lstlisting}

The list is short and stable; changes require explicit review.

\section{Adding a Seventh Boundary}
\label{sec:bridge_boundaries:adding}

The policy for adding a new bridge boundary:

\begin{enumerate}
  \item Document the user-visible need that the new boundary
    serves. If the need can be satisfied by routing through an
    existing boundary, do that instead.
  \item Document the external library that the new boundary will
    import.
  \item Add the file to the whitelist in
    \code{tools/check\_bridges.py}.
  \item Update this chapter (and the relevant API reference
    chapter) to include the new boundary.
\end{enumerate}

The policy has been invoked zero times since the rule was
established. Every additional integration since then has been
satisfied by extending one of the existing boundaries (e.g.,
adding safetensors support to \code{serialization/} or DLPack
support to \code{converters.py}). The intent is to keep the count
at six indefinitely.

\section{Summary and Forward References}
\label{sec:bridge_boundaries:summary}

\Lucid's engine is hermetically sealed against external numerical
libraries except at six explicitly named bridge boundaries. The
boundaries are: external array ingest (\code{converters.py}),
tensor-side egress (\code{tensor.py}), display
(\code{\_repr.py}), typing protocol (\code{\_types.py}),
checkpoint serialisation (\code{serialization/} and two optim
files), and external data ingest (\code{dataloader.py}). Hard Rule
H4 enforces the closure; a static check rejects any new import of
a forbidden library outside the whitelist.

The next chapter, \chref{ch:python_cpp_binding}, describes the
pybind11 layer that the Python side of \Lucid\ sits on.
\chref{ch:lazy_imports} closes \partref{part:system} with the
lazy-import surface that the top-level package presents to the
user. \chref{ch:serialization} returns to boundary 5 with a focus
on the serialisation formats themselves; \chref{ch:data_pipeline}
returns to boundary 6.

\begin{lucidnote}
For the user: the bridge boundaries are invisible from the API
surface. \code{lucid.from\_numpy(arr)} just works; \code{.numpy()}
just works; \code{lucid.save} and \code{lucid.load} just work.
The chapter is for the contributor who needs to understand where
numpy or DLPack can legitimately appear in the codebase and where
it cannot.
\end{lucidnote}
